\section{Paso 3: validación del modelo conceptual}
\label{sec:paso3-validacion-conceptual}
% Origen: 01-definicion-del-caso.md §7, 03-informe-tpi.md §4, 08-consulta-catedra.md, 10-antecedentes.md §2 y SIM.md (log del 20/09 y del 05/10).

La validación conceptual busca establecer que los supuestos del modelo representan el sistema con el
detalle que la pregunta exige, antes de programarlo \parencite{law2015}. Hubo tres instancias: la revisión
con la cátedra, el contraste con la literatura y la revisión de los supuestos contra los datos
disponibles.

\subsection{Revisión con la cátedra}

El tema fue aprobado por la cátedra el 20/09/2026, con el pedido de precisar la línea, la franja, los días
y las medidas de rendimiento. El caso definido en la sección~\ref{sec:paso1-formulacion}, con el argumento
de la congestión transitoria y la estructura en oleadas, se presentó en la reunión del 23/09. Después, la
cátedra aprobó el uso de la Pedestrian Library para representar el vestíbulo y pidió solicitar los planos de la
estación a los organismos responsables. Eso dio origen a los pedidos de información pública de la
sección~\ref{sec:paso2-datos-modelo-conceptual}.

Tres decisiones posteriores se tomaron de forma provisoria y se sometieron a la cátedra:
\begin{enumerate}
  \item redefinir \EI{} como los 22 molinetes del Principal en lugar de 28;
  \item trabajar con rangos con fuente para los parámetros sin dato publicado;
  \item dejar E2 y E3 como trabajo futuro.
\end{enumerate}
\pendiente{registrar la respuesta de la cátedra a estas tres decisiones.}

\subsection{Contraste con la literatura}

Cada supuesto de comportamiento se comparó con lo medido en estudios de estaciones de metro
(tabla~\ref{tab:contraste}).

\begin{table}[H]
\centering
\caption{Supuestos del modelo frente a la literatura}
\label{tab:contraste}
\small
\begin{tabularx}{\textwidth}{@{}>{\raggedright\arraybackslash}p{2.6cm} >{\raggedright\arraybackslash}p{3.2cm} X@{}}
\toprule
Supuesto & Valor del modelo & Referencia y lectura \\
\midrule
Tiempo de validación & Triangular (1,8; 2,4; 3,5)~s, media 2,57~s &
2~s en Valparaíso \parencite{seriani2025}; 2,91 a 3,49~s en Pekín \parencite{wang2024}. El valor
intermedio cae dentro de lo observado y el contexto pesimista, con media 3,23~s, cubre el caso de Pekín \\
Velocidad de marcha & Uniforme (1,1; 1,5)~m/s &
Rango de la literatura y valor calibrado de 1,37~m/s \parencite{li2024}. Las normas con valores menores
consideran pasajeros con equipaje, poco frecuentes en el pico de la mañana \parencite{tian2026} \\
Elección de molinete & Desvío lateral más espera estimada &
Fuerte sesgo por cercanía \parencite{seriani2025,tian2026} y peso de la fila \parencite{wang2024}.
Confirma que el término de distancia es necesario. Advierte también que la ubicación de los molinetes que
se agregan en \EI{} importa tanto como su cantidad \\
Arribos & Oleadas por horario del Roca & El antecedente con datos abiertos de molinetes usa arribos de
Poisson \parencite{zhen2024}. Las oleadas son la diferencia central de este modelo \\
\bottomrule
\end{tabularx}
\end{table}

\subsection{Revisión de los supuestos contra los datos}

\begin{itemize}
  \item \textbf{Cola por molinete, no cola única.} La primera versión lógica usaba una cola común con
        servidores intercambiables. Ese esquema no puede reproducir el desbalance observado entre
        molinetes (factor 11). Por eso el modelo principal usa una cola por molinete con elección
        explícita, que es lo que se ve en una batería de molinetes.
  \item \textbf{Volumen tomado de validaciones.} Usar las validaciones como demanda supone que los
        molinetes no se saturan. En el pico, cada molinete activo procesa en promedio unos 114 pasajeros
        por ventana, menos de la mitad de los 246 que llegó a registrar un molinete en una ventana. En
        promedio sobre 15 minutos no hay saturación. Sí puede haberla durante pocos minutos después de
        cada tren, y eso es lo que el modelo busca medir.
  \item \textbf{Cola vacía a las 07:00.} A las 07:00 el flujo es un 30~\% menor que en el pico, y entre
        trenes la capacidad sobra. Es razonable que la cola se vacíe entre oleadas al comienzo de la
        franja, pero no hay medición que lo confirme.
  \item \textbf{Molinetes de \EI{} disponibles.} Que los 22 identificadores correspondan a molinetes
        habilitables para ingreso es un supuesto. Se pidió la correspondencia a SBASE.
\end{itemize}

\subsection{Límites de la validación conceptual}

El grupo no pudo medir en Constitución ni consultar al personal operativo. Quedan sin confirmar:
\begin{itemize}
  \item el criterio real de habilitación de molinetes;
  \item la proporción de pasajeros del Roca que pasa a la Línea~C;
  \item la geometría del vestíbulo.
\end{itemize}
Por eso el modelo se considera validado conceptualmente en su estructura, que representa oleadas,
caminata, elección y cola, pero no en sus valores. Las conclusiones sobre la espera quedan condicionadas a
los rangos de la tabla~\ref{tab:parametros}.
